\chapter{Perfil Profissional}

O egresso do curso será um profissional especializado no desenvolvimento, implantação e manutenção de soluções baseadas em Inteligência Artificial, apto a atuar em empresas de tecnologia, indústrias, startups, centros de pesquisa, instituições públicas e organizações privadas.

Será capaz de integrar conhecimentos de Engenharia de Software, Ciência de Dados, Inteligência Artificial, Computação em Nuvem e Automação para desenvolver sistemas inteligentes que apoiem processos de tomada de decisão, automação industrial, desenvolvimento de software, processamento de linguagem natural, visão computacional e análise preditiva.

Ao concluir a formação, o profissional deverá demonstrar competência para:

\begin{itemize}
\item Projetar soluções de Inteligência Artificial.
\item Desenvolver aplicações baseadas em Machine Learning.
\item Desenvolver soluções utilizando Deep Learning.
\item Desenvolver aplicações utilizando IA Generativa.
\item Construir Agentes Inteligentes.
\item Desenvolver aplicações utilizando RAG.
\item Integrar APIs corporativas.
\item Desenvolver aplicações Full Stack assistidas por IA.
\item Implantar soluções em Cloud Computing.
\item Implantar modelos utilizando práticas de MLOps.
\item Desenvolver aplicações para IoT e Indústria 4.0.
\item Aplicar princípios de segurança da informação.
\item Atuar de forma ética e responsável no desenvolvimento de sistemas inteligentes.
\end{itemize}
